\documentclass[10pt,a4paper]{amsart}
\usepackage[margin=19mm]{geometry}
\usepackage{amsmath,amssymb,mathtools}
\usepackage[scr=rsfs,cal=euler]{mathalfa}
\usepackage{xfrac}
\usepackage[unicode,hidelinks,bookmarks=false]{hyperref}
\newcommand{\ie}{i.e.\ }
\newcommand{\cf}{cf.\ }
\newcommand{\resp}{respectively\ }
\newcommand{\Subsection}{\subsection}
\providecommand{\nobreakdash}{\nobreak\hbox{-}}
\setlength{\emergencystretch}{2em}
\multlinegap0pt
\usepackage{amssymb}
\newtheorem{lemma}{Lemma}
\title{A simple proof of exponential decay for the near-critical planar Ising model}
\author{Jianping Jiang, Frederik Ravn Klausen}
\date{Source: arXiv:2512.07554v2 (CC BY 4.0)}
\begin{document}
\maketitle

\noindent\emph{Source and licence.} A simple proof of exponential decay for the near-critical planar Ising model. Version arXiv:2512.07554v2 (CC BY 4.0), distributed by arXiv under the Creative Commons Attribution 4.0 International licence, http://creativecommons.org/licenses/by/4.0/, as recorded in the arXiv licence metadata for this exact version and shown on the abstract page https://arxiv.org/abs/2512.07554v2. Full licence notice: \texttt{licenses/arxiv-2512.07554-CC-BY-4.0.txt}.\par\smallskip

\noindent\textit{Editorial scope.} Counted unit: the article's lemma lem:moments (source0.tex lines 328--335). The article's complete original proof of it is delivered with it (source0.tex lines 336--364, one \texttt{proof} environment of 29 lines). No mathematics is written, repaired or re-derived: the statement and the proof are the article's own lines, in the article's own notation, and no retained line is substituted or re-flowed.  Retained uncounted as the source-local context the proof needs: lines 162--164; lines 165--167; lines 218--223; lines 279--284; lines 286--289.  One declared adaptation: exactly 1 ancillary strings inside the retained spans are omitted (image-inclusion commands and, where the source repeats a label, the repeated label definitions) and no other line differs from the source; the figure environments, with their placement, captions and labels, are kept, and the package records the omitted strings in fidelity receipts bound to the affected bodies.  No retained line contains a citation, so the wrapper carries no bibliography and no bibliography entry.  The retained lines contain the article's own diagram code, which the wrapper typesets with the standard tikz packages; no image file is delivered and the package declares no asset dependency.  Editorial formatting only, in addition: an AMS \texttt{amsart} preamble that restates the article's own declarations and macros used by the retained lines, namely the environments \texttt{lemma} on separate counters, together with the standard packages those lines need.  The retained environments print in this document as Lemma 1 in order of appearance, whereas the article prints the same items with the numbering of its own full document.  The complete native source of the article as distributed in the arXiv e-print for this exact version is bundled unchanged as \texttt{sources/190580/source0.tex}.
\begin{equation}\label{eq:nearcriticalRSW}
	\phi^1_{T,a,h}(\exists \text{ dual open circuit in }T^a\setminus S^a \text{ surrounding }S^a)>c_0,
\end{equation}

\begin{equation}\label{eq:onearm}
	C_1(a/r)^{1/8} \leq \phi^0_{\Lambda_{r},a,h}(0 \overset{a\mathbb{Z}^2}{\longleftrightarrow} \partial \Lambda^a_{r}) \leq \phi^1_{\Lambda_{r},a,h}(0 \overset{a\mathbb{Z}^2}{\longleftrightarrow} \partial \Lambda^a_{r}) \leq C_2 (a/r)^{1/8},~\forall r\in [a,1/2], \tag{1-arm}
\end{equation}

	\begin{figure}
		\begin{center}
			
			\caption{An illustration of the events $E(R)$ (left) and $H(R)$ (right). Here $T:=[0,10]\times[0,3]$, $R:=[2,8]\times [0,3]$, $S:=[1,9]\times[1,2]$, $Q_1:=[1,2)\times [1,2]$ and $Q_8:=(8,9]\times[1,2]$. The dotted circuit is a dual open circuit in $T^a\setminus S^a$. $S^a$ has a left-right crossing. Both $Q^a_1$ and $Q^a_8$ have a top-bottom crossing. In the crossing cluster of $S^a$ only the two marked points in $Q^a_1$ and $Q^a_8$ are connected to the ghost directly by one external edge. Note that $H(R)\subset E(R)$.}\label{fig:Looptog}
		\end{center}
	\end{figure}

		\begin{equation}\label{eq:E_1E_2}
			\begin{split}
				&E_1:=\{\exists \text{ dual open circuit in }T^a\setminus S^a \text{ surrounding }S^a\},\\
				&E_2:=\{S^a \text{ has LR open crossing}\}\cap \{Q^a_1 \text{ has TB open crossing}\} \cap \{Q^a_8 \text{ has TB open crossing}\},
			\end{split}
		\end{equation}

    \begin{equation}\label{eq:Ndef}
        \mathcal{N}:=|\mathscr{C}|\mathbf{1}[{E_2}], \ 
        \mathcal{N}_1:=|\mathscr{C}\cap Q_1^a|\mathbf{1}[{E_2}], \  \mathcal{N}_8:=|\mathscr{C}\cap Q_8^a|\mathbf{1}[{E_2}].
    \end{equation}

\begin{lemma}\label{lem:moments}
	See Figure \ref{fig:Looptog} for the definition of $S$ and $T$, \eqref{eq:E_1E_2} for $E_1$ and $E_2$, and \eqref{eq:Ndef} for $\mathcal{N},\mathcal{N}_1,\mathcal{N}_8$.
	There exists $h_0, C_5, C_6, C_7\in(0,\infty)$ such that for all $a\in(0,1]$ and $h\in[0,h_0]$,
	\begin{align*}
		&C_5 a^{-15/8}\leq \phi_{S,a,h}^0(\mathcal{N}_i)\leq C_6 a^{-15/8},~\phi_{S,a,h}^0(\mathcal{N}_i^2)\leq C_7 a^{-15/4} \text{ for }i=1,8.\\
		&\phi_{T,a,h}^{\xi}(\mathcal{N} \mid E_1)\leq C_6 a^{-15/8},~\forall \text{ boundary conditions }\xi.
	\end{align*}
\end{lemma}

\begin{proof}
    We may assume $a\in[0,1/100]$ since the $a\in(1/100,1]$ case follows from a finite energy argument. To simplify notation, all $\longleftrightarrow$ in the proof denote internal connection (i.e., using only edges in $a\mathbb{Z}^2$).
	Whenever $E_2$ occurs, each $x\in \mathscr{C}\cap Q^a_1$ must satisfy $x\longleftrightarrow \partial \Lambda_{1}^a(x)$. So, together with \eqref{eq:onearm},
	\[\phi_{S,a,h}^0(\mathcal{N}_1)\leq \phi_{S,a,h}^0\left(\sum_{x\in Q_1^a}\mathbf{1}[x\longleftrightarrow \partial \Lambda^a_{1}(x)]\mathbf{1}[E_2]\right)\leq \sum_{x\in Q_1^a}\phi_{S,a,h}^0(x\longleftrightarrow \partial \Lambda^a_{1}(x))\leq C_6 a^{-15/8}.\]
	
	The upper bound for $\phi_{T,a,h}^{\xi}(\mathcal{N} \mid E_1)$ follows in a similar way, using \eqref{eq:nearcriticalRSW}.
	
	To prove the lower bound, we denote by $E_3$ the event that there are enough LR and TB crossings of $Q^a_1$ such that any path in $Q_1^a$ with diameter $\geq 1/4$ hits one of those crossings. The critical RSW tells us that $\phi_{Q_1,a,0}^0(E_3)>0$. Let $D^a:=((1,1)+[1/4,3/4]^2)\cap a\mathbb{Z}^2$. Then by FKG and \eqref{eq:onearm}, 
	\begin{align*}
		\phi_{S,a,h}^0(\mathcal{N}_1) &\geq \phi_{S,a,h}^0(\mathcal{N}_1\mathbf{1}[E_3])\geq \phi_{S,a,h}^0\left(\sum_{x\in D^a}\mathbf{1}[x\longleftrightarrow\partial \Lambda_{1/4}^a(x)]\mathbf{1}[E_3]\mathbf{1}[E_2]\right)\\
		&\geq \phi_{S,a,h}^0(E_2)  \phi_{S,a,h}^0(E_3) \phi_{S,a,h}^0\left(\sum_{x\in D^a}\mathbf{1}[x\longleftrightarrow\partial \Lambda_{1/4}^a(x)]\right) \geq C_5 a^{-15/8}.
	\end{align*}

%So by linearity of expectation, FKG and \eqref{eq:onearm}, 
%   	\begin{align*}
%		\phi_{S,a,h}^0(\mathcal{N}_1) &\geq \phi_{S,a,h}^0(\mathcal{N}_1\mathbf{1}%[E_3])\geq  \sum_{x\in D^a} \phi_{S,a,h}^0\left(x\longleftrightarrow\partial %\Lambda_{1/4}^a(x),E_3\right) \geq C_5 a^{-15/8},
%	\end{align*}
%	where the third inequality follows from the FKG inequality, and the last inequality follows from \eqref{eq:onearm}.
	
	For the second moment, using $\{x \longleftrightarrow \mathscr{C}, y\longleftrightarrow \mathscr{C}, E_2\}\subset \{x \longleftrightarrow y\}$,

\begin{align*}\phi_{S,a,h}^0(\mathcal{N}_1^2)\leq \sum_{x,y\in Q_1^a} \phi_{S,a,h}^0(x\longleftrightarrow y)\leq |Q_1^a|+\sum_{x\in Q_1^a}\sum_{k=0}^{\lceil -\log_2 a \rceil}\sum_{y\in Q_1^a: 2^{k}a\leq |y-x|<2^{k+1}a} \phi_{S,a,h}^0(x\longleftrightarrow y).
	\end{align*}
	Note that for all $x,y\in Q_1^a$ with $|y-x|\geq 2^{k}a$, FKG and \eqref{eq:onearm} give
	\[\phi_{S,a,h}^0(x\longleftrightarrow y)\leq \phi_{\Lambda_{2^{k-2}a}(x),a,h}^1(x\longleftrightarrow \partial \Lambda_{2^{k-2}a}^a(x))\phi_{\Lambda_{2^{k-2}a}(y),a,h}^1(y\longleftrightarrow \partial \Lambda_{2^{k-2}a}^a(y))\leq[\tilde{C}_2 2^{-(k-2)/8}]^2,\]
	  where we changed $C_2$ to $\tilde{C}_2$ to include the $k\in\{0,1\}$ case. Plugging this into the last displayed inequality, we get
	\[\phi_{S,a,h}^0(\mathcal{N}_1^2)\leq C_7a^{-15/4}.\]
The proofs for the inequalities involving $Q_8$ are the same.	 
\end{proof}

\end{document}
